\chapter{Fazit und Ausblick}

\section{Fazit}
Ziel der Entwicklung des GDB-Stubs für D3OS war die Möglichkeit, das Betriebssystem auf echter Hardware debuggen zu können. Anfangs bot das System nur die Möglkichkeit in einer emulierten Umgebung einen Debugger zu verwenden. Durch die Implementierung des Stubs kann der Kernel von D3OS, und in begrenztem Ausmaß auch der User-Space, auf echter Hardware debuggt werden. Der Stub ermöglicht die wichtigsten Funktionen eines Debuggers, wie das Setzen und Entfernen von Breakpoints, auslesen von Speicher und die schrittweise Ausführung des Programms. Wie in der Evaluation aber bereits beschrieben, dürfen aufgrund der Funktionsweise des GDB RSP und der Exception Handler keine Breakpoint oder Step-Into Befehle für bestimmte Funktionsaufrufe verwendet werden. Um bestimmte Funktionen vom Step-Into auszuschließen, lässt sich der \texttt{skip} Befehl verwenden.
Außerdem ist der Stub noch nicht kompatibel mit Multi-Core-Scheduling und funktioniert auf echter Hardware nur auf einer älteren Version von D3OS aufgrund der in der Evaluation beschriebenen Probleme bei der Datenübertragung. Im Abgaberepository \footnote{https://git.hhu.de/bsinfo/thesis/ba-mof81wap} unter der branch \texttt{abgabe} zu dieser Arbeit befindet sich die Version des Stubs, die auf Hardware funktioniert und einen alten Stand von D3OS nutzt. Die Entwicklungsversion, die nur auf QEMU zuverlässig funktioniert, aber die neuste Version von D3OS unterstützt, ist unter der branch \texttt{development} zu finden.

\section{Ausblick}
Die Implementierung lässt sich in folgenden Bereichen weiterentwickeln:
\begin{itemize}
\item \textbf{User-Space Kompatibilität:} Aktuell lässt sich der User-Space ausschließlich über Hardware Breakpoints debuggen. Für alle anderen Befehle muss der Adressraum geladen werden, was der Stub noch nicht macht. Um Nutzer-Anwendungen richtig zu debuggen, muss der Stub hier erweitert werden.
\item \textbf{Neuste Version von D3OS:} Der Stub funktioniert in QEMU mit der neusten Version von D3OS und dem neuen Multi-Core-Scheduler. Dazu muss in der Konfiguration ein Core eingestellt werden. Auf echter Hardware treten aber noch die in der Evaluation beschriebenen Probleme auf.  Der Stub kann hier um Multi-Core-Kompatibilität erweitert werden. Die Probleme des seriellen Ports können ebenfalls weiter untersucht werden.
\item \textbf{TCP statt Serial:} Alternativ zum seriellen Port kann TCP verwendet werden. Dazu benötigt D3OS Treiber, die eine TCP-Verbindung bei der Verwendung auf echter Hardware zulässt. Sobald eine Verbindung hergestellt werden kann, kann die Implementierung der \texttt{GdbConnection} auf TCP geändert werden. Die Implementierung der einzelnen Schnittstellen des Targets und der EventLoop müssen dazu nicht verändert werden.
\item \textbf{Weitere GDB Feature:} Das GDB RSP und die Crate \texttt{gdbstub} bieten noch viele Erweiterungen für verschiedene Anwendungsfälle an. Durch die Nutzung von IDETs der Crate lässt sich die Implementierung leicht um weitere Traits erweitern. Eine sinnvolle Erweiterung wären zum Beispiel Hardware Watchpoints. Dadurch lassen sich Breakpoints auf Speicherzugriffe setzen.
\end{itemize}
